\documentclass[11pt]{article}

\usepackage{amsmath}
\usepackage{setspace}
\usepackage{amssymb}
\usepackage{changepage}
\usepackage{graphicx}
\usepackage[a4paper, margin=0.5in]{geometry}

\DeclareMathOperator*{\argmin}{arg\,min}

\setstretch{1.3}

\hbadness=10000
\setlength{\parindent}{0pt}

\newcommand\numberthis{\addtocounter{equation}{1}\tag{\theequation}}

\begin{document}
    \begin{center}
        \Large \textbf{STAT3006 Worksheet 06}
    \end{center}

    \Large \textbf{1. Read AdaBoost}

    \large Real AdaBoost is an additive boosting method that constructs a strong classifier by 
    iteratively
    adding weak learners with real-valued outputs. Unlike discrete AdaBoost, each weak learner
    contributes a confidence-weighted prediction. \\ 

    The additive model takes the form 

    \begin{align*}
        \begin{matrix}
            h^{(B)}(\mathbf{x}) = \sum_{b = 1}^B \nu f^{(b)}(\mathbf{x}),
            & f^(b)(\mathbf{x}) \in \mathbb{R}
        \end{matrix}
    \end{align*}

    where $\nu > 0$ is a shrinkage parameter. In this exercise, we set $\nu = 1$ so that 

    \begin{align*}
        h^{(b)}(\mathbf{x}) = \sum_{b = 1}^B f^(b)(\mathbf{x})
    \end{align*}

    The final classifier is 

    \begin{align*}
        \hat{y} = \text{sign}\bigg(h^{(b)}(\mathbf{x})\bigg)
    \end{align*}

    Real AdaBoost is derived by minimising the exponential loss

    \begin{align*}
        \begin{matrix}
        L(h) = \sum_{i = 1}^n \exp\bigg(-y_i h(x_i)\bigg),
        & y_i \in \{-1, +1 \}
        \end{matrix}
    \end{align*}

    (a) Show that minimising the exponential loss after adding a weak learner 

    \begin{align*}
        h^{(b)}(\mathbf{x}) = h^{(b - 1)}(\mathbf{x}) + f^{(b)}(\mathbf{x})
    \end{align*}

    is equivalent to solving a weighted exponential loss problem: 

    \begin{align*}
        \min\limits_{j} \sum_{i = 1}^n w^{(b)}_i \exp \bigg(-y_i f(\mathbf{x}_i)\bigg)
    \end{align*}

    where the weights are 

    \begin{align*}
        w^{(b)}_i = \exp \bigg(-y_i h^{(b - 1)}(\mathbf{x}_i)\bigg)
    \end{align*}

    \textbf{Answer:} The loss function is given by

    \begin{align*}
        L(h) 
        &= \sum_{i = 1}^n \exp \bigg(-y_i h(\mathbf{x}_i)\bigg) \\ 
        &= \sum_{i = 1}^n \exp \bigg\{-y_i \bigg(h^{(b - 1)}(\mathbf{x}_i) + f^{(b)}(\mathbf{x}_i)
            \bigg)\bigg\} \\ 
        &= \sum_{i = 1}^n \exp \bigg\{-y_i h^{(b - 1)}(\mathbf{x}_i)\bigg\} \cdot \exp \bigg\{-y_i
            f^{(b)}(\mathbf{x}_i)\bigg\}
    \end{align*}

    Result follows \\ 

    (b) Assume a weak learner is a tree with Regions $R_1, \dots, R_j$:
    
    \begin{align*}
        \begin{matrix}
            f^{(b)}(\mathbf{x}) = \gamma_j, 
            & \text{ if } \mathbf{x} \in R_j
        \end{matrix}
    \end{align*}

    Show that the optimal leaf value in Real AdaBoost is 

    \begin{align*}
        \gamma_j = \frac{1}{2}\log \frac{\sum_{i \in R_j, y_i = 1} w^{(b)}_i}{\sum_{i \in R_j, y_i =
        -1}w^{(b)}_i}
    \end{align*}

    \textbf{Answer:} First we define the loss function for a given region $j$, 

    \begin{align*}
        L_j(h) = \sum_{i = 1}^n w^{(b)}_i \exp \bigg\{-y_i \gamma_j\bigg\}
    \end{align*}

    We partition $L_j$ into two summations: one where $y_i = -1$ and another one where $y_i = +1$,

    \begin{align*}
        L_j(h) = \sum_{i \in R_j, y_i = 1} w^{(b)}_i \exp (-\gamma_j) + \sum_{i \in R_j, y_1 = -1}
        w^{(b)}_i \exp (\gamma_j)
    \end{align*}

    Differentiating with respect to $\gamma_j$ we get

    \begin{align*}
        \frac{\partial L_j(h)}{\partial \gamma_j}
        &= -\sum_{i \in R_j, y_i = 1} w^{(b)}_i \exp(-\gamma_j) + \sum_{i \in R_j, y_i = -1}
        w^{(b)}_i \exp(\gamma_j)
    \end{align*}

    Letting L.H.S. = 0 we get 

    \begin{align*}
        \sum_{i \in R_j, y_i = 1} w^{(b)}_i \exp (-\gamma_j) 
        =& \sum_{i \in R_j, y_i = -1} w^{(b)}_i \exp (\gamma_j) \\ 
        \frac{\exp(\gamma_j)}{\exp(\gamma_j)}
         &= \frac{\sum_{i \in R_j, y_i = 1}w^{(b)}_i}{\sum_{i \in R_j, y_i = -1}w^{(b)}_i} \\ 
         \exp(2 \gamma_j)
         &= \frac{\sum_{i \in R_j, y_i = 1}w^{(b)}_i}{\sum_{i \in R_j, y_i = -1}w^{(b)}_i} \\ 
         2\gamma_j
         &= \log \frac{\sum_{i \in R_j, y_i = 1}w^{(b)}_i}{\sum_{i \in R_j, y_i = -1}w^{(b)}_i} \\ 
         \gamma_j 
         &= \frac{1}{2} \log \frac{\sum_{i \in R_j, y_i = 1}w^{(b)}_i}{\sum_{i \in R_j, y_i = -1}w^{(b)}_i} \\ 
    \end{align*}

    (c) In standard boosting, one often writes

    \begin{align*}
        h^{(b)}(\mathbf{x}) = h^{(b - 1)}(\mathbf{x}) + \alpha^{(b)}f^{(b)}(\mathbf{x}) \numberthis
    \end{align*}

    Explain where the step size $\alpha^{(b)}$ appears in Real AdaBoost \\

    \textbf{Answer:} In standard AdaBoost, also called discrete AdaBoost, since $y_i \in \{-1,
    +1,\}$ we have a step size $\alpha^{(b)}$ which transforms the terms in the summation from
    integer values to real values. \\ 

    In Real AdaBoost, the predictions aren't discrete and involve both the direction as well as
    confidence of the prediction. That is, $y \in [0, 1]$. Therefore there isn't an explicit
    $\alpha^{(b)}$ parameter and is otherwise encoded into the predictor function. \\ 

    (d) Explain why Real AdaBoost is more flexible than Discrete AdaBoost \\ 

    \textbf{Answer:} In Discrete AdaBoost, the $\alpha^{(b)}_i$ term is imposed onto all 
    predictions from that learner. Whereas with Real AdaBoost, this term is encoded straight into
    the predictor as each prediction has both direction and confidence. We can think of it as every
    single prediction in every learner has an $\alpha^{(b)}$ term of its own, making Real AdaBoost
    more flexible.
\end{document}
